% Generated draft: requires mathematical and source review.
\documentclass[11pt]{article}
\usepackage[margin=25mm]{geometry}
\usepackage{fontspec}
\setmainfont{Latin Modern Roman}
\usepackage{amsmath,amssymb,mathtools,unicode-math}
\setmathfont{Latin Modern Math}
\usepackage{enumitem,needspace,xurl,hyperref}
\setlength{\parindent}{0pt}
\setlength{\parskip}{5pt}
\setlength{\emergencystretch}{3em}
\newcommand{\N}{\mathbb{N}}
\newcommand{\Z}{\mathbb{Z}}
\newcommand{\Q}{\mathbb{Q}}
\newcommand{\R}{\mathbb{R}}
\newcommand{\C}{\mathbb{C}}
\newcommand{\F}{\mathbb{F}}
\newcommand{\A}{\mathbb{A}}
\newcommand{\PP}{\mathbb{P}}
\newcommand{\E}{\mathbb{E}}
\newcommand{\eps}{\varepsilon}
\newcommand{\dd}{\,\mathrm{d}}
\begin{document}
\section*{UnsolvedMath ID: 30005712}
\textbf{Source problem code:} catalogue entry
\subsection{Definitions and mathematical statement}
The supplied record asks whether every simplicial sphere is generically anisotropic over every field, but does not define generic anisotropy. The following is an explicit Stanley--Reisner interpretation requiring verification against the cited report.

Let \(k\) be a field and let \(D\) be a finite simplicial complex on vertices \(\{1,\ldots,n\}\) triangulating \(S^{d-1}\), where \(d\geq 1\). Define its Stanley--Reisner ideal and face ring by
\[
 I_D=\left(\prod_{j\in F}x_j : F\subseteq\{1,\ldots,n\},\ F\notin D\right),
 \qquad k[D]=k[x_1,\ldots,x_n]/I_D,
\]
with \(\deg x_j=1\). Let \(t_{ij}\), for \(1\leq i\leq d\) and \(1\leq j\leq n\), be algebraically independent over \(k\), and put
\[
 K=k(t_{ij}:1\leq i\leq d,\ 1\leq j\leq n),\qquad
 \theta_i=\sum_{j=1}^{n}t_{ij}x_j,\qquad
 A_K(D)=K[D]/(\theta_1,\ldots,\theta_d).
\]
Under the interpretation used here, generic anisotropy over \(k\) means that this generic Artinian reduction satisfies
\[
 \text{for every integer }q\geq0\text{ with }2q\leq d,
 \quad u\in A_K(D)_q\setminus\{0\}\ \Longrightarrow\ u^2\neq0
 \text{ in }A_K(D)_{2q}.
\]
The main question is whether this property holds for every field \(k\), every integer \(d\geq1\), and every such simplicial sphere \(D\). The use of a purely transcendental coefficient extension makes the interpretation meaningful also for finite fields. Neither this precise construction nor the intended convention for simplicial spheres is specified in the supplied excerpt; their agreement with the source remains unresolved.

The original statement also supplies contextual implications, labelled ``Theorem 6 ([3])'': for any field \(k\) and simplicial sphere \(D\), generic anisotropy of the suspension \(S(D)\) implies that \(E[D]\) has the Weak Lefschetz property; generic anisotropy of both \(D\) and \(S(D)\) implies that \(E[D]\) has the Strong Lefschetz property. Here \(S(D)=D*S^0\), where \(S^0\) consists of two isolated vertices. For a finite-dimensional standard graded algebra \(B\) over a field, the Weak Lefschetz property means that some \(\ell\in B_1\) makes every map \(\ell:B_i\to B_{i+1}\) have maximal rank. The Strong Lefschetz property means that some \(\ell\in B_1\) makes every map \(\ell^r:B_i\to B_{i+r}\), for \(i,r\geq0\), have maximal rank. The excerpt does not define \(E\) or \(E[D]\), so these supplied implications cannot be identified unconditionally with the reduction defined above; their reference [3] is also unresolved.

The excerpt additionally contains ``Question 8'': identify classes of Gorenstein standard graded algebras having the generic anisotropy property. A standard graded algebra over a field \(F\) is a graded algebra \(B=\bigoplus_{i\geq0}B_i\) with \(B_0=F\), generated by finitely many elements of \(B_1\). The excerpt does not specify the generic construction, coefficient fields, dimension restrictions, or degree range intended for this broader question, so it is retained as a separate, incompletely specified target.
\subsection{Short English statement}
The main question asks whether the generic Artinian face-ring reduction of every simplicial sphere has no nonzero homogeneous element whose square vanishes in degrees at most half the socle degree, over every field. This interpretation requires source verification. The supplied excerpt also records Lefschetz implications and a broader question about Gorenstein standard graded algebras, with missing notation and hypotheses.
\subsection{Sources}
\begin{itemize}
\item Local UnsolvedMath registry, permanent ID \(30005712\), source problem code \texttt{catalogue entry}, title ``Generic Anisotropy of Simplicial Spheres.'' The supplied record is the basis of this statement-only draft. It names no proposer and records proposed year 2023.
\item Supplied primary citation: \emph{Geometric, Algebraic and Topological Combinatorics} (2024), Oberwolfach Reports, pp.\ 3249--3318, DOI \texttt{10.4171/owr/2023/58}: \url{https://doi.org/10.4171/owr/2023/58}. The record describes these reports as documenting mathematical workshops and their open-problem discussions. The citation uses 2024, while the DOI and proposed-year metadata use 2023; these details are retained as supplied.
\item Supplied literature source: Papadakis--Petrotou, ``characteristic 2 anisotropy'': \url{https://arxiv.org/abs/2012.09815}.
\item Supplied literature source: Fan, ``strongly edge decomposable spheres'': \url{https://arxiv.org/abs/2406.18989}.
\item Supplied literature source labelled ``original Oberwolfach report'': \url{https://doi.org/10.4171/owr/2023/58}.
\item Literature assessment dated exactly \texttt{2026-08-22}: the supplied curation labels the problem partially solved and states that generic anisotropy is proved in characteristic 2 and for several sphere classes and low dimensions, while the assertion for every simplicial sphere over every field remains open. This is supplied dated curation derived from the OWR record, not an independently verified literature assessment; its status and cited claims require independent verification.
\item The registry labels the statement \texttt{corrected\_verified} and reports that it was reconstructed and checked against the official Oberwolfach report. This describes supplied verification metadata, not verification performed for this draft. The excerpt's internal reference [3] has no supplied bibliographic identification.
\end{itemize}
\end{document}
